% ====================================================================
% BAB 5: KESIMPULAN DAN SARAN
% ====================================================================
% File ini berisi kesimpulan dan saran penelitian yang mencakup:
% - Kesimpulan (menjawab rumusan masalah secara komprehensif)
% - Saran (rekomendasi akademik dan praktis)
% ====================================================================

\chapter{KESIMPULAN DAN SARAN}

% --- BAGIAN 5.1: KESIMPULAN ---
\section{Kesimpulan}
\label{sec:conclusion}

Berdasarkan seluruh rangkaian tahapan penelitian, perancangan sistem, implementasi, serta pengujian dan evaluasi kuantitatif yang telah dilaksanakan pada data percakapan grup WhatsApp, diperoleh beberapa kesimpulan utama yang menjawab rumusan masalah penelitian sebagai berikut:

\begin{enumerate}
    \item \textbf{Penerapan Modified BERTopic pada Percakapan WhatsApp:} \\
    Penelitian ini berhasil merancang dan menerapkan modifikasi arsitektur \textit{framework} BERTopic yang disesuaikan secara spesifik untuk menangani karakteristik wacana percakapan grup WhatsApp berbahasa Indonesia. Integrasi alur \textit{Single-Path Preprocessing} berhasil mereduksi derau (\textit{noise}) dan membersihkan 91.868 pesan mentah menjadi 56.089 dokumen berkualitas tinggi. Pada tahap pemodelan, pemanfaatan model bahasa \textit{embedding} IndoBERTweet (vektor 768 dimensi), reduksi dimensi UMAP (5 komponen dengan metrik kosinus), algoritma pengklasteran hierarkis BIRCH (\textit{BirchWithOutliers} dengan \textit{threshold} 0,25 dan batas persentil pencilan 95,0), representasi BM25 ($k_1 = 1,2; b = 0,5$) yang diperkaya kamus \textit{STOPWORDS} 9 kategori, serta reduksi topik otomatis dengan mekanisme pengaman (\textit{safety guard}) terbukti mampu mengekstraksi topik percakapan secara terstruktur dan adaptif.

    \item \textbf{Kualitas dan Performa Model Usulan Dibandingkan Baseline:} \\
    Hasil pengujian kurva pembelajaran (\textit{learning curve}) pada proporsi 20\% hingga 100\% data menunjukkan bahwa model usulan (\textit{proposed}) secara konsisten mengungguli model \textit{baseline} BERTopic standar pada seluruh aspek evaluasi kuantitatif:
    \begin{enumerate}[label=\alph*.]
        \item \textit{Pengurangan Outlier Rate}: Penggunaan algoritma BIRCH berhasil menekan persentase dokumen pencilan secara drastis dari 42,42\% pada \textit{baseline} menjadi hanya \textbf{8,79\%} pada model usulan (pada 100\% data), sehingga mempertahankan lebih dari 91\% volume percakapan ke dalam klaster topik yang valid.
        \item \textit{Peningkatan Koherensi Topik (C-NPMI)}: Model usulan menghasilkan nilai koherensi semantik kata kunci yang meningkat signifikan dari 0,0662 (\textit{baseline}) menjadi \textbf{0,1567}, membuktikan bahwa representasi IndoBERTweet dan BM25 mampu menghasilkan kata kunci topik yang padu dan relevan secara makna percakapan informal.
        \item \textit{Stabilitas Keberagaman Topik (Topic Diversity)}: Nilai \textit{Topic Diversity} pada model usulan terjaga sangat tinggi pada skor \textbf{0,9559} (rentang 0,9335--0,9919), yang mengonfirmasi bahwa kata kunci antartopik tidak mengalami tumpang-tindih (\textit{redundancy}) berkat penyaringan \textit{stopwords} 9 kategori.
        \item \textit{Kepadatan Vektor dan Kemiripan Internal Klaster}: Model usulan mencapai \textit{Embedding Density} sebesar \textbf{0,6666} (vs 0,5611 pada \textit{baseline}) dan \textit{Intra-topic Similarity} sebesar \textbf{0,6716} (vs 0,5211 pada \textit{baseline}), menegaskan bahwa klaster topik yang terbentuk memiliki batas kerapatan yang tegas serta tingkat homogenitas semantik yang erat.
    \end{enumerate}

    \item \textbf{Perancangan dan Implementasi Sistem Web Terintegrasi:} \\
    Penelitian ini berhasil merancang dan mengimplementasikan sistem analisis percakapan grup WhatsApp berbasis web (\textit{ChatAnalisis}) dengan arsitektur \textit{client-server} yang memisahkan antarmuka pengguna (\textit{frontend} Vue.js 3, TypeScript, TailwindCSS, Pinia) dan layanan pemrosesan komputasi (\textit{backend} FastAPI terakselerasi GPU/CUDA). Sistem mengintegrasikan fitur privasi data melalui \textit{client-side salted SHA-256 anonymization}, pemrosesan latar belakang asinkron (\textit{BackgroundTasks}), pemantauan progres \textit{real-time} berbasis \textit{Server-Sent Events} (SSE), penelusuran konteks pesan historis (\textit{message context traceability}), proteksi keamanan \textit{path traversal}, dan pembersihan berkas disk otomatis. Pengujian fungsional sistem menggunakan metode \textit{Black-Box Testing} terhadap 34 skenario kasus uji menunjukkan bahwa seluruh fitur pada modul antarmuka dan API \textit{backend} beroperasi 100\% valid sesuai dengan spesifikasi kebutuhan.
\end{enumerate}

Secara keseluruhan, integrasi IndoBERTweet, BIRCH \textit{clustering}, dan skema representasi BM25 pada \textit{framework} BERTopic yang diwujudkan dalam aplikasi web interaktif telah terbukti secara empiris dan aplikatif mampu menjadi solusi komprehensif dalam mengatasi tantangan \textit{information overload} serta ekstraksi wawasan tematik pada percakapan grup WhatsApp berskala besar.

% --- BAGIAN 5.2: SARAN ---
\section{Saran}
\label{sec:recommendations}

Berdasarkan hasil temuan penelitian, evaluasi sistem, serta keterbatasan yang dihadapi selama pelaksanaan penelitian, berikut adalah beberapa saran yang dapat diajukan untuk pengembangan akademik selanjutnya dan implementasi praktis:

\subsection{Saran untuk Penelitian Selanjutnya}
\label{subsec:future-research}

\begin{enumerate}
    \item \textbf{Eksplorasi Model Bahasa dan Arsitektur Embedding Mutakhir:} \\
    Penelitian selanjutnya dapat mengeksplorasi perbandingan performa menggunakan model bahasa generasi terbaru, seperti model berbasis \textit{Large Language Models} (LLM) berbahasa Indonesia atau model representasi kalimat dwibahasa mutakhir, guna menguji peningkatan pemahaman semantik pada percakapan dengan tingkat \textit{code-mixing} yang lebih kompleks.

    \item \textbf{Pengembangan Pemodelan Topik Dinamis (Dynamic Topic Modeling):} \\
    Penelitian dapat diperluas dengan menambahkan dimensi waktu (\textit{temporal dynamics}) untuk mengamati pergeseran dan evolusi topik percakapan dari waktu ke waktu (\textit{topic drift}), sehingga tren pembicaraan komunitas dapat dipetakan secara kronologis per periode tertentu.

    \item \textbf{Integrasi Analisis Wacana Lanjutan:} \\
    Pengembangan dapat mengombinasikan pemodelan topik dengan analisis sentimen, analisis emosi, atau analisis jejaring sosial (\textit{Social Network Analysis}) untuk memahami polaritas opini anggota kelompok dan pola interaksi sosial di dalam setiap tema pembicaraan.

    \item \textbf{Pengujian pada Domain dan Format Perangkat Beragam:} \\
    Disarankan untuk menguji generalisabilitas model pada domain percakapan grup yang berbeda (seperti grup koordinasi kerja, komunitas akademik, atau layanan publik) serta memperluas format data masukan dari sistem operasi lain (misalnya berkas ekspor chat dari perangkat iOS).
\end{enumerate}

\subsection{Saran untuk Praktisi dan Pengembangan Terapan}
\label{subsec:practical-recommendations}

\begin{enumerate}
    \item \textbf{Dukungan Multi-Platform Aplikasi Pesan Instan:} \\
    Pengembang aplikasi dapat menambahkan fitur penguraian (\textit{parser}) untuk platform pesan instan lainnya, seperti Telegram, Discord, atau Slack, sehingga platform web analisis ini dapat digunakan secara universal oleh berbagai komunitas digital.

    \item \textbf{Peningkatan Skalabilitas dan Manajemen Antrean Server:} \\
    Untuk implementasi pada lingkungan produksi berskala besar dengan banyak pengguna bersamaan (\textit{concurrent users}), disarankan untuk mengintegrasikan sistem antrean tugas terdistribusi (seperti Celery dan Redis) serta orkestrasi kontainer (Docker/Kubernetes).

    \item \textbf{Fitur Ekspor Laporan dan Notifikasi Berkala:} \\
    Disarankan untuk menambahkan fitur ekspor laporan hasil analisis topik ke dalam format PDF/Excel interaktif serta modul ringkasan berkala yang dapat dikirimkan secara otomatis kepada administrator grup.
\end{enumerate}
